% Part: lambda-calculus
% Chapter: syntax
% Section: substitution

\documentclass[../../../include/open-logic-section]{subfiles}

\begin{document}

\olfileid[fr]{lam}{syn}{sub}
\olsection{Substitution}

\begin{explain}
Les variables libres renvoient à des variables
de l'environnement. Il est donc naturel de
remplacer une variable libre par une valeur
particulière. Par exemple, on peut vouloir
remplacer $f$ dans $\lambd[x][f x]$ par
un terme donné, tel que la fonction identité
$\lambd[y][y]$. On obtient alors
$\lambd[x][(\lambd[y][y]) x]$. Cette opération,
qui remplace des variables libres par des
lambda-termes, s'appelle \emph{substitution}.
\end{explain}

\begin{defn}[Substitution] \ollabel{defn:substitution}
La \emph{substitution} d'un terme $N$ à une
variable $x$ dans un terme $M$,
$\Subst{M}{N}{x}$, est définie inductivement
par :
\begin{enumerate}
  \item $\Subst{x}{N}{x} = N$.
    \ollabel{defn:substitution-1}
  \item $\Subst{y}{N}{x} = y$ si $x \neq y$.
    \ollabel{defn:substitution-2}
  \item $\Subst{PQ}{N}{x} =
    (\Subst{P}{N}{x})(\Subst{Q}{N}{x})$.
    \ollabel{def:substitution-3}
  \item $\Subst{(\lambd[y][P])}{N}{x} =
    \lambd[y][\Subst{P}{N}{x}]$ si
    $x \neq y$ et $y \notin \FV{N}$;
    dans les autres cas, la substitution
    n'est pas définie.
    \ollabel{defn:substitution-4}
\end{enumerate}
\end{defn}

\begin{explain}
Dans la \olref{defn:substitution}\olref{defn:substitution-4},
la condition $x \neq y$ empêche de remplacer
par~$N$ les occurrences \emph{liées} de la
variable~$x$ dans~$M$. Par exemple, si l'on
considère la substitution
$\Subst{\lambd[x][x]}{y}{x}$, son résultat
devrait être $\lambd[x][x]$, et non
$\lambd[x][y]$.

Quand on substitue $N$ à~$x$ dans
$\lambd[y][P]$, on exige aussi que
$y \notin \FV{N}$. Par exemple, on ne peut
pas substituer $y$ à $x$ dans
$\lambd[y][x]$, c'est-à-dire former
$\Subst{\lambd[y][x]}{y}{x}$ selon la règle
ci-dessus : on obtiendrait
$\lambd[y][y]$, terme qui désigne la
fonction renvoyant directement l'argument
qu'elle reçoit. Or $\lambd[y][x]$ désigne
une fonction qui renvoie toujours le terme~$x$
(ou la valeur désignée par $x$). Le résultat
recherché serait donc une fonction qui reçoit
un argument, l'ignore et renvoie la variable
d'environnement~$y$. Pour construire ce résultat
correctement, il faudrait d'abord « renommer »
la variable liée~$y$.
\end{explain}

\begin{explain}
La définition donne une opération partielle. Lorsque le lieur
porte le nom~$x$, la quatrième clause ne s'applique pas, même
si l'opération usuelle sur les occurrences libres laisserait
l'abstraction intacte. Lorsqu'un lieur risque de capturer une
variable libre de~$N$, il faut d'abord renommer ce lieur.
La section suivante expliquera comment choisir un représentant
pour lequel la substitution est définie.
\end{explain}

\begin{prob}
Quels sont les résultats des substitutions
suivantes, lorsqu'elles sont définies?
\begin{enumerate}
  \item $\Subst{\lambd[y][x(\lambd[w][vwx])]}{(uv)}{x}$
  \item $\Subst{\lambd[y][x(\lambd[x][x])]}{(\lambd[y][xy])}{x}$
  \item $\Subst{y(\lambd[v][xv])}{(\lambd[y][vy])}{x}$
\end{enumerate}
\end{prob}

\begin{thm} \ollabel{thm:notinfv}
Si $x \notin \FV{M}$, alors
$\FV{\Subst{M}{N}{x}} = \FV{M}$,
pourvu que la substitution soit définie.
\end{thm}

\begin{proof}
Raisonnons par induction sur la formation de $M$.
\begin{enumerate}
\item $M$ est une variable : exercice.
\item $M$ est de la forme $(PQ)$ : exercice.
\item $M$ est de la forme $\lambd[y][P]$.
  Puisque $\Subst{\lambd[y][P]}{N}{x}$ est
  définie, elle vaut
  $\lambd[y][\Subst{P}{N}{x}]$.
  La substitution $\Subst{P}{N}{x}$ est donc
  définie; de plus, $x \neq y$ et
  $x \notin \FV{P}$. Ainsi :
  \begin{align*}
    \FV{\Subst{\lambd[y][P]}{N}{x}}
    &= \FV{\lambd[y][\Subst{P}{N}{x}]}
       && \text{par \olref{defn:substitution-4}}\\
    &= \FV{\Subst{P}{N}{x}} \setminus \{y\}
       && \text{par la \olref[fv]{def:fv}\olref[fv]{def:fv2}}\\
    &= \FV{P} \setminus \{y\}
       && \text{par l'hypothèse d'induction}\\
    &= \FV{\lambd[y][P]}
       && \text{par la \olref[fv]{def:fv}\olref[fv]{def:fv2}}.
  \end{align*}
\end{enumerate}
\end{proof}

\begin{explain}
Dans le cas de l'abstraction, $x$ et $y$ sont distinctes.
L'absence de~$x$ parmi les variables libres de l'abstraction
entraîne alors son absence parmi celles du corps~$P$.
\end{explain}

\begin{prob}
Compléter la preuve du
\olref[lam][syn][sub]{thm:notinfv}.
\end{prob}

\begin{thm} \ollabel{thm:infv}
Si $x \in \FV{M}$, alors
$\FV{\Subst{M}{N}{x}} =
(\FV{M} \setminus \{x\}) \cup \FV{N}$,
pourvu que la substitution soit définie.
\end{thm}

\begin{proof}
Raisonnons par induction sur la formation de $M$.
\begin{enumerate}
\item $M$ est une variable : exercice.
\item $M$ est de la forme $PQ$. Puisque
  $\Subst{(PQ)}{N}{x}$ est définie, elle vaut
  $(\Subst{P}{N}{x})(\Subst{Q}{N}{x})$,
  les deux substitutions étant définies.
  Comme $x \in \FV{PQ}$, on a soit
  $x \in \FV{P}$, soit $x \in \FV{Q}$,
  soit les deux. La suite est laissée
  en exercice.
\item $M$ est de la forme $\lambd[y][P]$.
  Puisque $\Subst{\lambd[y][P]}{N}{x}$ est
  définie, elle vaut
  $\lambd[y][\Subst{P}{N}{x}]$,
  avec $\Subst{P}{N}{x}$ définie,
  $x \neq y$ et $y \notin \FV{N}$.
  Puisque $x \in \FV{\lambd[y][P]}$,
  on a aussi $x \in \FV{P}$. Dès lors :
  \begin{align*}
    \FV{\Subst{(\lambd[y][P])}{N}{x}}
    &= \FV{\lambd[y][\Subst{P}{N}{x}]}\\
    &= \FV{\Subst{P}{N}{x}} \setminus \{y\}\\
    &= ((\FV{P} \setminus \{x\}) \cup \FV{N})
       \setminus \{y\}
       && \text{par l'hypothèse d'induction}\\
    &= (\FV{P} \setminus \{x,y\}) \cup \FV{N}
       && \text{car }y \notin \FV{N}\\
    &= (\FV{\lambd[y][P]} \setminus \{x\})
       \cup \FV{N}.
  \end{align*}
\end{enumerate}
\end{proof}

\begin{explain}
Dans le calcul sous l'abstraction, $x$ est retirée des variables
libres du corps lors de la substitution. Le lieur retire
ensuite~$y$. Comme $y\notin\FV{N}$, il ne retire aucune des
variables libres apportées par~$N$.
\end{explain}

\begin{prob}
Compléter la preuve du
\olref[lam][syn][sub]{thm:infv}.
\end{prob}

\begin{thm}\ollabel{thm:clr}
Si $x \notin \FV{N}$ et si
$\Subst{M}{N}{x}$ est définie, alors
$x \notin \FV{\Subst{M}{N}{x}}$.
\end{thm}

\begin{proof}
Exercice.
\end{proof}

\begin{prob}
Démontrer le \olref[lam][syn][sub]{thm:clr}.
\end{prob}

\begin{thm}\ollabel{thm:inv}
Si $\Subst{M}{y}{x}$ est définie et
$y \notin \FV{M}$, alors
$\Subst{\Subst{M}{y}{x}}{x}{y} = M$.
\end{thm}

\begin{proof}
Raisonnons par induction sur la formation de $M$.
\begin{enumerate}
\item $M$ est une variable $z$ : exercice.
\item $M$ est de la forme $(PQ)$. Alors :
  \begin{align*}
    \Subst{\Subst{(PQ)}{y}{x}}{x}{y}
    &=\Subst{((\Subst{P}{y}{x})(\Subst{Q}{y}{x}))}{x}{y} \\
    &= (\Subst{\Subst{P}{y}{x}}{x}{y})
       (\Subst{\Subst{Q}{y}{x}}{x}{y}) \\
    &= (PQ) && \text{par l'hypothèse d'induction}.
  \end{align*}
\item $M$ est de la forme $\lambd[z][N]$.
  Puisque $\Subst{\lambd[z][N]}{y}{x}$
  est définie, on a $z \neq x$ et $z \neq y$.
  De plus, $y \notin \FV{M}$ entraîne
  $y \notin \FV{N}$. L'hypothèse d'induction
  assure que la substitution inverse sur~$N$
  est définie. Ainsi :
  \begin{align*}
    \Subst{\Subst{(\lambd[z][N])}{y}{x}}{x}{y}
    &= \Subst{(\lambd[z][\Subst{N}{y}{x}])}{x}{y} \\
    &= \lambd[z][\Subst{\Subst{N}{y}{x}}{x}{y}] \\
    &= \lambd[z][N]
       && \text{par l'hypothèse d'induction}.
  \end{align*}
\end{enumerate}
\end{proof}

\begin{prob}
Compléter la preuve du
\olref[lam][syn][sub]{thm:inv}.
\end{prob}

\end{document}
